% Self-contained source. Compile with: latexmk -pdf erdos_727_partial_proof.tex
% No external bibliography, figures, or nonstandard local files are required.
% Author and AI-assistance identifiers supplied by Johan Land.
\documentclass[11pt,a4paper]{article}
\usepackage[T1]{fontenc}
\usepackage[utf8]{inputenc}
\usepackage{newtxtext}
\usepackage{amsmath,amssymb,amsthm,mathtools}
\usepackage{newtxmath}
\usepackage[margin=26mm,headheight=14pt,headsep=7mm]{geometry}
\usepackage{microtype,booktabs,array,needspace}
\usepackage[font=small,labelfont=bf]{caption}
\usepackage{fancyhdr}
\usepackage[hidelinks]{hyperref}
\hypersetup{pdftitle={A carry-counting approach to the k=3 case of Erdos Problem 727},
  pdfsubject={Candidate proof; draft for review},pdfauthor={Johan Land}}

\newcommand{\paperauthor}{Johan Land}
\newcommand{\paperaffiliation}{Independent Researcher}
\newcommand{\N}{\mathbb N}
\newcommand{\Z}{\mathbb Z}
\newcommand{\R}{\mathbb R}
\newcommand{\LL}{\mathcal L}
\newcommand{\1}{\mathbf 1}
\DeclareMathOperator{\Li}{Li}
\newtheorem{theorem}{Theorem}[section]
\newtheorem{lemma}[theorem]{Lemma}
\newtheorem{corollary}[theorem]{Corollary}
\numberwithin{equation}{section}

\pagestyle{fancy}
\fancyhf{}
\fancyhead[L]{\small\itshape Erd\H{o}s Problem 727: the case \(k=3\)}
\fancyhead[R]{\small Draft for review}
\fancyfoot[C]{\small\thepage}
\renewcommand{\headrulewidth}{0.3pt}
\setlength{\parindent}{1.2em}
\setlength{\parskip}{0pt}
\setlength{\emergencystretch}{1em}
\setlength{\textfloatsep}{12pt plus 3pt minus 3pt}
\setlength{\floatsep}{10pt plus 2pt minus 2pt}
\renewcommand{\arraystretch}{1.18}
\allowdisplaybreaks[1]
\raggedbottom

\title{\vspace{-1.3em}A carry-counting approach to the \(k=3\) case\\
of Erd\H{o}s Problem 727}
\author{\paperauthor\\[0.25em]\small\paperaffiliation}
\date{6 September 2026}

\begin{document}
\maketitle
\vspace{-1.2em}
\begin{abstract}
We present a candidate unconditional proof that infinitely many positive integers
\(n\) satisfy \(((n+3)!)^2\mid(2n)!\), and hence also \(((n+2)!)^2\mid(2n)!\).
The quadratic \(f(t)=210t^2+391t+179\) makes each of \(f(t)+1,f(t)+2,f(t)+3\)
a product of two linear forms. A fixed arithmetic progression handles small
primes; direct carry counting, uniform quadratic exponential-sum estimates,
and a finite residue calculation bound the remaining failures. The only
external analytic input is the classical prime number theorem in fixed
arithmetic progressions. The argument gives a positive proportion of successful
parameters on the progression and \(\gg\sqrt Z\) successful integers \(n\leq Z\).
\end{abstract}
\vspace{0.2em}
\noindent\textbf{Status.} This is a candidate proof circulated for scrutiny. The accompanying
Lean development is complete on the author's machine (see the closing section);
independent expert verification and independent reproduction of that build
remain pending. The general case of arbitrary fixed \(k\) is not
claimed here.

\section{Statement and strategy}

Erd\H{o}s, Graham, Ruzsa and Straus asked whether, for every fixed \(k\geq2\),
there are infinitely many \(n\) for which \(((n+k)!)^2\mid(2n)!\)
\cite[p.~90]{EGRS}. We consider the case \(k=3\), using
\[
 f(t)=210t^2+391t+179.
\]

\begin{theorem}\label{thm:main}
There exist fixed integers \(Q\geq1\), \(0\leq t_0<Q\), and \(X_0\geq1\)
such that for every integer \(X\geq X_0\), at least \(X/100\) integers
\(v\in[X,2X)\) satisfy
\begin{equation}\label{eq:target}
 \bigl((f(t_0+Qv)+3)!\bigr)^2\mid\bigl(2f(t_0+Qv)\bigr)!.
\end{equation}
Consequently there is a constant \(c_*>0\) such that, for all sufficiently
large \(Z\),
\[
 \#\{1\leq n\leq Z:((n+3)!)^2\mid(2n)!\}\geq c_*\sqrt Z.
\]
\end{theorem}

\begin{corollary}\label{cor:k2}
There are infinitely many positive integers \(n\) such that
\(((n+2)!)^2\mid(2n)!\).
\end{corollary}

Write \(A(n)=(n+1)(n+2)(n+3)\) and \(B(n)=\binom{2n}{n}\).
Since \((n+3)!=n!A(n)\) and \((2n)!=(n!)^2B(n)\), cancellation gives
\begin{equation}\label{eq:cancel}
 ((n+3)!)^2\mid(2n)!\quad\Longleftrightarrow\quad A(n)^2\mid B(n).
\end{equation}
For a prime \(p\), Legendre's formula yields the carry identity
\begin{equation}\label{eq:carry}
 c_p(n):=v_p(B(n))
 =\sum_{h\geq1}\left(\left\lfloor\frac{2n}{p^h}\right\rfloor
       -2\left\lfloor\frac n{p^h}\right\rfloor\right)
 =\sum_{h\geq1}\1_{\{n/p^h\}\geq1/2}.
\end{equation}
Thus it suffices to prove \(c_p(n)\geq2v_p(A(n))\) for every prime dividing
\(A(n)\). We arrange this for all bounded primes by congruences, then bound
the union of the remaining failure events. No independence between primes
or between the six linear factors is assumed.

Here \(\{x\}=x-\lfloor x\rfloor\), \(e(x)=\exp(2\pi i x)\), and \(\|x\|\)
denotes distance to the nearest integer. All sums indexed by \(p\) are over
primes. All moduli, constants, and exponent margins are fixed before
\(X\to\infty\); implicit constants may depend on them.

\section{The polynomial family and a fixed progression}\label{sec:arithmetic}

The three factorizations are
\begin{equation}\label{eq:factors}
\begin{aligned}
 f(t)+1&=(35t+36)(6t+5),\\
 f(t)+2&=(t+1)(210t+181),\\
 f(t)+3&=(15t+14)(14t+13).
\end{aligned}
\end{equation}
Let \(L_i(t)=a_it+b_i\), \(1\leq i\leq6\), denote the displayed factors,
and let \(j_i\in\{1,2,3\}\) be the associated shift. Table~\ref{tab:forms}
records all constants needed below. Direct expansion verifies
\begin{equation}\label{eq:coordinates}
 f(t)=\alpha_iL_i(t)^2+\beta_iL_i(t)-j_i,
 \qquad f'(-b_i/a_i)=c_i.
\end{equation}
\begin{table}[ht]
\centering
\caption{The six factors; \(d_i\) is the reduced denominator of \(\alpha_i\).}
\label{tab:forms}
\begin{tabular}{@{}clrrrrr@{}}
\toprule
\(i\)&\(L_i(t)\)&\(j_i\)&\(c_i\)&\(\alpha_i\)&\(\beta_i\)&\(d_i\)\\
\midrule
1&\(35t+36\)&1&\(-41\)&\(6/35\)&\(-41/35\)&35\\
2&\(6t+5\)&1&41&\(35/6\)&\(41/6\)&6\\
3&\(t+1\)&2&\(-29\)&210&\(-29\)&1\\
4&\(210t+181\)&2&29&\(1/210\)&\(29/210\)&210\\
5&\(15t+14\)&3&\(-1\)&\(14/15\)&\(-1/15\)&15\\
6&\(14t+13\)&3&1&\(15/14\)&\(1/14\)&14\\
\bottomrule
\end{tabular}
\end{table}

Every \(L_i\) is primitive, and all prime factors of the slopes lie in
\(\{2,3,5,7\}\). The within-pair resultants have absolute values
\(41,29,1\). A prime dividing factors from different pairs divides a
nonzero difference between two of \(f+1,f+2,f+3\), so is at most~2.
Consequently a prime \(p>1000\) divides at most one \(L_i(t)\), and its
root class satisfies \(f'(t)\equiv c_i\not\equiv0\pmod p\).

\subsection{Small primes by lifting and the Chinese remainder theorem}

Set \(Y=1000\) and \(G(t)=\prod_iL_i(t)=A(f(t))\). Since
\(G(0)=180\cdot181\cdot182\), its prime support is
\(S=\{2,3,5,7,13,181\}\). For every prime \(p\leq Y\) outside \(S\),
impose \(t\equiv0\pmod p\), ensuring \(p\nmid G(t)\).
For \(p\in S\), set \(e_p=v_p(G(0))\) and choose a lifting depth
\(\lambda_p\) as follows:
\begin{center}
\begin{tabular}{@{}lrrrrrr@{}}
\toprule
\(p\)&2&3&5&7&13&181\\
\midrule
\(e_p\)&3&2&1&1&1&1\\
\(\lambda_p\)&3&3&2&2&2&2\\
\(\lambda_p+2e_p\)&9&7&4&4&4&4\\
\bottomrule
\end{tabular}
\end{center}
Each \(\lambda_p\) is larger than every individual valuation \(v_p(179+j)\),
\(1\leq j\leq3\). If \(r_p\) is the least nonnegative residue of~179
modulo \(p^{\lambda_p}\), prescribe
\begin{equation}\label{eq:local}
 f(t)\equiv r_p+p^{\lambda_p}(p^{2e_p}-1)\pmod{p^{\lambda_p+2e_p}},
 \qquad t\equiv0\pmod p.
\end{equation}
This has a unique lift of the root \(t=0\pmod p\), because
\(f'(0)=391=17\cdot23\) is coprime to every \(p\in S\).
Indeed, for \(r\geq1\),
\[
 f(u+zp^r)\equiv f(u)+zp^rf'(u)\pmod{p^{r+1}};
\]
when the derivative is a unit, varying \(z\pmod p\) realizes each next
target digit exactly once.

Condition~\eqref{eq:local} preserves each \(v_p(f(t)+j)=v_p(179+j)\)
and prescribes \(2e_p\) consecutive digits equal to \(p-1\) in positions
\(\lambda_p,\ldots,\lambda_p+2e_p-1\) of \(f(t)\). Each produces a carry on doubling,
even for \(p=2\), regardless of the incoming carry. Hence
\(c_p(f(t))\geq2e_p=2v_p(G(t))\).
The Chinese remainder theorem now gives \(0\leq t_0<Q\), with
\begin{equation}\label{eq:Q}
 Q=\prod_{p\in S}p^{\lambda_p+2e_p}
   \prod_{\substack{p\leq Y\\p\notin S}}p,
\end{equation}
such that every \(t=t_0+Qv\) satisfies the required inequalities for all
\(p\leq Y\). The modulus is fixed, however large it may be.

Write
\[
 F(v)=f(t_0+Qv),\qquad \LL_i(v)=L_i(t_0+Qv).
\]
For \(v\in[X,2X)\), all six forms are positive and at most \(CX\), for
some fixed \(C\geq1\), and \(F(v)\asymp X^2\). For \(p>Y\), the slope
\(a_iQ\) is a unit modulo every power of \(p\).

\subsection{Repeated factors and failure events}

Discard parameters for which \(p^2\mid\LL_i(v)\) for some \(i\) and
\(p>Y\). Each condition is one class modulo \(p^2\), so their union has size
\begin{equation}\label{eq:square}
 E_{\rm sq}(X)\leq6X\sum_{p>Y}\frac1{p^2}
       +6\pi(\sqrt{CX})
 \leq\frac{6X}{Y}+o(X)=\frac3{500}X+o(X).
\end{equation}
Here \(\sum_{m>Y}m^{-2}\leq1/Y\), and \(\pi(z)\leq z\) suffices for
the endpoint error. For a remaining parameter, every prime \(p>Y\)
dividing \(G(t)\) has exponent exactly one. If \(p\mid\LL_i(v)\), then
\(F(v)\equiv-j_i\pmod p\); since \(p>6\), this supplies a carry at
level~1. It remains to exclude
\begin{equation}\label{eq:bad}
 p\mid\LL_i(v),\qquad
 \{F(v)/p^h\}<\tfrac12\quad\text{for every }h\geq2.
\end{equation}
We bound the union of these events in the ranges below.

\section{Prime harmonic sums}\label{sec:primes}

The only external analytic theorem used here is the classical estimate
\begin{equation}\label{eq:pnt}
 \pi(x;d,a)=\frac{\Li(x)}{\varphi(d)}
       +O_d\bigl(x\exp(-c_d\sqrt{\log x})\bigr),
 \qquad (a,d)=1,
\end{equation}
for fixed \(d\), with \(c_d>0\); see \cite[\S27.12, (27.12.5) and
(27.12.8)]{DLMF}. Here \(\pi(x;d,a)\) counts primes in the indicated
class, \(\varphi\) is Euler's totient, and
\(\Li(x)=\int_2^x du/\log u\). In particular,
\(\pi(CX)=o(X)\) for fixed \(C\).

\begin{lemma}\label{lem:prime}
For fixed \(d,a\) with \((a,d)=1\),
\begin{equation}\label{eq:mertens}
 \sum_{\substack{p\leq z\\p\equiv a\ (d)}}\frac1p
 =\frac{\log\log z}{\varphi(d)}+B_{d,a}+o(1).
\end{equation}
Consequently, if \(R\) is a fixed set of reduced classes modulo \(d\),
\(0<\alpha<\beta\), and \(C>0\) is fixed, then
\begin{equation}\label{eq:bandprimes}
 \sum_{\substack{X^\alpha<p\leq CX^\beta\\p\bmod d\in R}}\frac1p
 =\frac{|R|}{\varphi(d)}\log\frac\beta\alpha+o(1).
\end{equation}
Also, for fixed \(P>1\), \(c>0\), and \(\delta>0\),
\begin{equation}\label{eq:weighted}
 \lim_{X\to\infty}\sum_{P<p\leq X^\delta}\frac1p
       \exp\!\left(-\frac{c\log X}{\log p}\right)
 =\int_0^\delta e^{-c/y}\frac{dy}{y}
 =E_1(c/\delta),
\end{equation}
where \(E_1(z)=\int_z^\infty e^{-w}\,dw/w\).
\end{lemma}

\begin{proof}
Partial summation writes the left side of~\eqref{eq:mertens} as
\(\pi(z;d,a)/z+\int_2^z\pi(u;d,a)u^{-2}\,du\). The main term
in~\eqref{eq:pnt} gives \((\log\log z)/\varphi(d)\) plus a constant;
the error integral converges because
\(\int_2^\infty e^{-c_d\sqrt{\log u}}\,du/u<\infty\).
Taking differences gives~\eqref{eq:bandprimes}, including the unrestricted
prime sum when \(d=1\).

For~\eqref{eq:weighted}, write the unrestricted sum as
\(H(z)=\log\log z+B+R_0(z)\), where \(R_0(z)\to0\), and integrate
\(g_X(z)=\exp(-c\log X/\log z)\) against \(dH(z)\) on \((P,X^\delta]\).
The main term is
\(\int_{\log P/\log X}^{\delta}e^{-c/y}\,dy/y\); its missing lower tail
tends to zero. To control the error, fix \(Z>P\). The part below~\(Z\)
tends to zero as \(X\to\infty\). On \([Z,X^\delta]\), integration by
parts and monotonicity of \(g_X\) bound the error by
\(2\sup_{z\geq Z}|R_0(z)|g_X(X^\delta)\).
First let \(X\to\infty\), then \(Z\to\infty\). The substitution \(w=c/y\)
gives the last equality in~\eqref{eq:weighted}.
\end{proof}

\section{Very small primes}\label{sec:small}

Fix
\[
 \delta=\frac1{20},\qquad \rho=\frac{1001}{2000},\qquad
 c=\frac12\log(1/\rho).
\]
For \(Y<p\leq X^\delta\), put
\(\ell=\lfloor\log X/(2\log p)\rfloor\). Then \(\ell\geq10\) and
\(p^\ell\leq\sqrt X\). On the unique root class of \(\LL_i\pmod p\),
\(F'(v)\equiv Qc_i\not\equiv0\pmod p\). The lifting argument above
therefore makes \(v\mapsto F(v)\) a bijection from that class modulo
\(p^\ell\) onto the residues congruent to \(-j_i\pmod p\).

The units digit is \(p-j_i\) and already gives a carry. To avoid a carry
at the next level, there are \((p-1)/2\) choices for the next digit;
with no incoming carry thereafter, each later digit has \((p+1)/2\)
choices. Thus the number of residues avoiding every carry at levels
\(2,\ldots,\ell\) is
\[
 \frac{p-1}{2}\left(\frac{p+1}{2}\right)^{\ell-2}
 \leq p^{\ell-1}\rho^{\ell-1}.
\]
Each residue occurs at most \(X/p^\ell+1\) times, giving
\begin{equation}\label{eq:smalllocal}
 \#\{v\in[X,2X):\text{\eqref{eq:bad} holds for }i,p\}
 \leq\left(\frac Xp+p^{\ell-1}\right)\rho^{\ell-1}.
\end{equation}
The sum of the second terms, over all primes in this range and all six
forms, is \(O(\sqrt X\log\log X)=o(X)\). Since
\[
 \ell-1\geq\frac{\log X}{2\log p}-2,
 \qquad \rho^{\ell-1}\leq\rho^{-2}
                  \exp\!\left(-\frac{c\log X}{\log p}\right),
\]
Lemma~\ref{lem:prime} bounds the union count by
\begin{equation}\label{eq:small}
 \limsup_{X\to\infty}\frac{E_{\rm small}(X)}X
 \leq6\rho^{-2}E_1(c/\delta)<\frac1{100}.
\end{equation}
For the strict numerical bound, \(\rho^{-2}<4\),
\(c/\delta=10\log(2000/1001)>6\), and \(E_1(z)\leq e^{-z}/z\);
hence the expression is less than \(4e^{-6}<1/100\).
The comparisons follow, without numerical integration, from
\(\log x\geq2(x-1)/(x+1)\) for \(x\geq1\) and
\(e^3>\sum_{r=0}^8 3^r/r!>20\).

\section{Medium primes and uniformity}\label{sec:medium}

\subsection{A quadratic exponential-sum bound}

\begin{lemma}\label{lem:quadratic}
Let \(q\geq3\) be odd, \((a,q)=1\), and let \(I\) be an interval of
\(L\) consecutive integers. Uniformly in these data and in \(\theta\in\R\),
\begin{equation}\label{eq:quadratic}
 \left|\sum_{r\in I}e\!\left(\frac{ar^2}{q}+\theta r\right)\right|
 \ll\left(\frac L{\sqrt q}+\sqrt q\right)\log(2q).
\end{equation}
\end{lemma}

\begin{proof}
Suppose first \(L\leq q\), and translate \(I\) to \(0\leq r<L\).
For \(b\pmod q\), set
\(G_b=\sum_{s\bmod q}e((as^2-bs)/q)\). Squaring and putting the
difference of the two variables equal to \(h\) gives
\[
 |G_b|^2=\sum_{h\bmod q}e\!\left(\frac{ah^2-bh}{q}\right)
                 \sum_{s\bmod q}e\!\left(\frac{2ahs}{q}\right)=q,
\]
since \(2a\) is a unit modulo \(q\). Fourier inversion yields
\[
 \sum_{r=0}^{L-1}e\!\left(\frac{ar^2}{q}+\theta r\right)
 =\frac1q\sum_{b\bmod q}G_b
                  \sum_{r=0}^{L-1}e\bigl((\theta+b/q)r\bigr).
\]
The inner absolute value is at most
\(\min(L,(2\|\theta+b/q\|)^{-1})\), interpreted as \(L\) at integral
frequencies. The points \(\theta+b/q\pmod1\) are \(1/q\)-spaced, so
only a bounded number lie in each distance shell of width \(1/q\) from
an integer. Summing the geometric-sum bounds gives \(O(q\log(2q))\),
uniformly in \(\theta\). Since \(|G_b|=\sqrt q\), the desired short-interval
bound is \(O(\sqrt q\log(2q))\). For general \(L\), divide \(I\) into
at most \(L/q+1\) intervals of length at most~\(q\).
\end{proof}

\subsection{Uniform distribution in each band}

Fix \(\eta=10^{-6}\) and \(\gamma_j=2/j\), \(4\leq j\leq40\).
For \(J=4,\ldots,39\), use the band
\begin{equation}\label{eq:bands}
 X^{\gamma_{J+1}+\eta}<p\leq X^{\gamma_J-\eta}.
\end{equation}
These exponent intervals are nonempty, since
\(\gamma_J-\gamma_{J+1}\geq1/780>2\eta\).

Fix \(i\) and choose the least nonnegative root \(v_0\) of
\(\LL_i\pmod p\). Conditional on \(p\mid\LL_i(v)\), write
\(v=v_0+pz\), where \(z\) runs over \(L=X/p+O(1)\) consecutive
integers. With \(D=210Q^2\),
\begin{equation}\label{eq:expandedphase}
 F(v_0+pz)=Dp^2z^2+pF'(v_0)z+F(v_0).
\end{equation}
We show that the vectors
\[
 \bigl(\{F(v)/p^2\},\ldots,\{F(v)/p^J\}\bigr)
\]
are asymptotically uniformly distributed in \([0,1)^{J-1}\),
\emph{uniformly over primes in the band}.

Fix a nonzero integer frequency \((h_2,\ldots,h_J)\), and let \(H\)
be its highest nonzero index. If \(H\geq3\), the quadratic coefficient
of the phase \(\sum_{r=2}^Jh_rF(v_0+pz)/p^r\) has reduced denominator
\(q=p^{H-2}\), for all sufficiently large \(X\). Indeed, its numerator
modulo \(p\) is \(Dh_H\neq0\pmod p\) once \(p\) exceeds the prime
factors of the fixed nonzero integer \(Dh_H\).
Lemma~\ref{lem:quadratic}, allowing the arbitrary real linear coefficient,
gives the normalized bound
\begin{equation}\label{eq:uniformquad}
 O\!\left(\left(p^{-(H-2)/2}+\frac{p^{H/2}}X\right)\log X\right)=o(1)
\end{equation}
uniformly in~\eqref{eq:bands}: here \(p\geq X^\delta\), \(H\leq J\), and
\(p^{H/2}/X\leq p^{J/2}/X\leq X^{-\eta J/2}\).
If \(H=2\), the quadratic coefficient is the integer \(Dh_2\), and the
nonconstant phase reduces to \(h_2F'(v_0)z/p\). Because
\(F'(v_0)\equiv Qc_i\pmod p\), its numerator modulo \(p\) is the fixed
nonzero integer \(h_2Qc_i\). The geometric sum is \(O(p)\), with
normalized bound
\begin{equation}\label{eq:uniformlinear}
 O(p/L)=O(p^2/X)=o(1),
\end{equation}
uniformly, since \(p\leq X^{1/2-\eta}\).

For completeness, fixed-frequency cancellation suffices uniformly for the
box \([0,1/2)^{J-1}\). Approximate its indicator above and below by
continuous periodic functions whose integrals differ arbitrarily little,
then approximate those functions by finite trigonometric polynomials.
At each stage only finitely many frequencies occur, so
\eqref{eq:uniformquad}--\eqref{eq:uniformlinear} hold uniformly for all
of them. First let \(X\to\infty\), then refine the approximation.
It follows that, uniformly over the band, the proportion of conditioned
parameters without a carry at levels \(2,\ldots,J\) is
\begin{equation}\label{eq:box}
 2^{1-J}+o(1).
\end{equation}

Summing~\eqref{eq:box} times \(X/p+O(1)\) over the primes incurs only
\(o(X)\): the reciprocal-prime sum in a fixed band is bounded by
Lemma~\ref{lem:prime}, the discrepancy is uniform, and the endpoint
terms total \(O(\pi(X^{1/2-\eta}))=o(X)\). Summing the six forms and
all 36 bands therefore gives
\begin{equation}\label{eq:medium}
\begin{aligned}
 \limsup_{X\to\infty}\frac{E_{\rm med}(X)}X
 &\leq6\sum_{J=4}^{39}2^{1-J}\log\frac{J+1}{J}\\
 &\leq6\cdot\frac14\sum_{J=4}^{\infty}2^{1-J}
 =\frac38.
\end{aligned}
\end{equation}
We used \(\log(1+1/J)\leq1/J\leq1/4\).

\subsection{Boundary strips}

The prime exponents between \(\delta=1/20\) and \(1/2+\eta\) not covered
by~\eqref{eq:bands} lie in the 37 strips
\[
 X^{\gamma_j-\eta}\leq p\leq X^{\gamma_j+\eta},\qquad 4\leq j\leq40.
\]
Count every event \(p\mid\LL_i(v)\) in these strips, ignoring carries.
Lemma~\ref{lem:prime} and \(X/p+O(1)\) give
\begin{equation}\label{eq:strips}
\begin{aligned}
 \limsup_{X\to\infty}\frac{E_{\rm bdry}(X)}X
 &\leq6\sum_{j=4}^{40}\log\frac{\gamma_j+\eta}{\gamma_j-\eta}\\
 &\leq12\eta\sum_{j=4}^{40}\frac1{\gamma_j-\eta}
 \leq\frac{6\eta\cdot814}{1-20\eta}<\frac1{100}.
\end{aligned}
\end{equation}
Here \(\sum_{j=4}^{40}j=814\); endpoint errors are \(o(X)\), since all
primes involved are at most \(X^{1/2+\eta}\). Any overlaps at boundaries
are harmless in these union bounds.

\section{Large primes and a finite residue calculation}\label{sec:large}

Suppose \(p>X^{1/2+\eta}\) divides \(L_i(t)\), with \(t=t_0+Qv\), and
write \(L_i(t)=pm\). Then \(m\geq1\), and uniformly in this range
\begin{equation}\label{eq:correction}
 \frac mp\leq\frac{CX}{p^2}=O(X^{-2\eta})=o(1),\qquad
 \frac{f(t)}{p^2}=\alpha_i m^2+\beta_i\frac mp-\frac{j_i}{p^2}.
\end{equation}
For \(i\neq3\), Table~\ref{tab:forms} shows that \(d_i\mid a_i\) and
\((b_i,d_i)=1\). Reducing \(L_i(t)=pm\) modulo \(d_i\) gives
\begin{equation}\label{eq:mclass}
 m\equiv b_ip^{-1}\pmod{d_i}.
\end{equation}
As the unit class of \(p\) varies, this permutes the unit classes of
\(m\). The relation uses the original coefficients \(a_i,b_i\), so
restricting \(t\) to the CRT progression does not change it.

The possible fractional parts of \(\alpha_i m^2\), for unit \(m\), are
listed in Table~\ref{tab:residues}. Each distinct value has equally many
preimages: squaring is a homomorphism on the finite unit group with
constant fiber size, and the numerator of \(\alpha_i\) is a unit modulo
\(d_i\). Thus the last column is the fraction of unit classes for which
this fractional part is below \(1/2\).
\begin{table}[ht]
\centering
\caption{Large-prime residues. Divide the listed numerators by \(d_i\)
to obtain the fractional parts.}
\label{tab:residues}
\begin{tabular}{@{}crlc@{}}
\toprule
\(i\)&\(d_i\)&Numerators of the fractional parts&\(\theta_i\)\\
\midrule
1&35&\(6,19,24,26,31,34\)&\(1/6\)\\
2&6&\(5\)&0\\
4&210&\(1,79,109,121,151,169\)&\(1/3\)\\
5&15&\(11,14\)&0\\
6&14&\(1,9,11\)&\(1/3\)\\
\bottomrule
\end{tabular}
\end{table}

No listed fraction is \(0\) or \(1/2\). The finite sets have positive
distance from both the integers and \(1/2\). By~\eqref{eq:correction},
the correction tends uniformly to zero, so for sufficiently large \(X\)
it neither crosses \(1/2\) nor wraps around an integer. Every listed
fraction above \(1/2\) therefore supplies a carry at level \(p^2\).
Failures lie in a fixed set \(R_i\) of reduced prime classes modulo
\(d_i\), with \(|R_i|/\varphi(d_i)=\theta_i\), by~\eqref{eq:mclass}.

The form \(i=3\) is favorable for a different reason:
\[
 \frac{f(t)}{p^2}=210m^2-\frac{29m}{p}-\frac2{p^2}.
\]
Its correction is strictly negative and tends uniformly to zero. It
lies in \((-1/2,0)\) for large \(X\), so the fractional part lies in
\((1/2,1)\). This always gives the second carry; put \(\theta_3=0\).

Since \(p\leq CX\) whenever \(p\mid\LL_i(v)\), the large-prime union count
satisfies
\[
 E_{\rm large}(X)\leq
 \sum_{i\neq3}\ \sum_{\substack{X^{1/2+\eta}<p\leq CX\\
                          p\bmod d_i\in R_i}}
                      \left(\frac Xp+O(1)\right).
\]
The total endpoint error is \(O(\pi(CX))=o(X)\), even though \(C\) may
be large, because it is fixed. Lemma~\ref{lem:prime} and
\(\sum_i\theta_i=1/6+1/3+1/3=5/6\) give
\begin{equation}\label{eq:large}
 \limsup_{X\to\infty}\frac{E_{\rm large}(X)}X
 \leq\frac56\log\frac1{1/2+\eta}
 <\frac56\log2<\frac7{12}.
\end{equation}
The last inequality follows from \(\log2<7/10\); the first four terms
of the exponential series already give \(e^{7/10}>2\).

\section{Completion of the proof}\label{sec:completion}

\begin{proof}[Proof of Theorem~\ref{thm:main} and Corollary~\ref{cor:k2}]
All primes \(p\leq Y\) satisfy the carry inequalities throughout the
fixed progression. After the repeated-factor exclusions, every failure
at a larger prime lies in~\eqref{eq:bad}. The very small range, medium
bands, boundary strips, and large range cover every such prime, and no
prime divisor of any \(\LL_i(v)\) exceeds \(CX\).

Let \(E(X)\) be the number of unsuccessful parameters in \([X,2X)\).
The union bound, together with \eqref{eq:square}, \eqref{eq:small},
\eqref{eq:medium}, \eqref{eq:strips}, and \eqref{eq:large}, gives
\begin{equation}\label{eq:budget}
 \limsup_{X\to\infty}\frac{E(X)}X
 <\frac3{500}+\frac1{100}+\frac38+\frac1{100}+\frac7{12}
 =\frac{2953}{3000}<\frac{99}{100}.
\end{equation}
All constants and dimensions are fixed; the Fourier approximations were
fixed before taking each asymptotic limit and only then refined. The
finitely many estimates therefore hold simultaneously on a common tail.
For every sufficiently large integer \(X\), at least \(X/100\) of the
\(X\) integers in \([X,2X)\) succeed. By~\eqref{eq:cancel} this is
exactly~\eqref{eq:target}.

The map \(v\mapsto F(v)\) is strictly increasing for positive \(v\), so
these parameters give infinitely many distinct successful integers.
Also \(F(v)\leq C_2v^2\) for some fixed \(C_2>0\) and all sufficiently
large~\(v\). Given large \(Z\), take
\[
 X=\left\lfloor\frac12\sqrt{Z/C_2}\right\rfloor.
\]
Every \(v\in[X,2X)\) then has \(F(v)\leq Z\), and at least
\(X/100\gg\sqrt Z\) of these distinct values succeed. This proves the
counting assertion. Finally, \(((n+2)!)^2\mid((n+3)!)^2\), so every
successful \(n\) for \(k=3\) is also successful for \(k=2\).
\end{proof}

\section*{Relation to the accompanying Lean development}

This note condenses the original candidate manuscript of 6 September 2026,
retaining its analytic route and its five-term failure budget. The accompanying
Lean source follows a reorganized proof: explicit local witnesses replace
iterated lifting, proved Mertens estimates in fixed progressions replace the
stronger quoted prime number theorem, and finite Fourier analysis replaces
the limiting box approximation. Its failure budget also differs: the boundary
strips are absorbed into the medium range, and the four range bounds are
slackened to \(0.01+0.02+0.35+0.60=0.98<1\), which suffices for infinitude.
It should not be described as a line-by-line certification of this
presentation.

The source endpoints \texttt{Erdos727.erdos727\_k3} and
\texttt{Erdos727.erdos727\_k2} express the intended infinitude statements.
The former unfolds to
\begin{quote}\small\ttfamily
Set.Infinite \{n : \(\N\) |\\
\hspace*{1em}(Nat.factorial (n + 3))\,\^{}\,2\ \(\mid\)\ Nat.factorial (2 * n)\}
\end{quote}
The stronger quantitative conclusions above are not separately exported
as named theorems in that development. At commit \texttt{c0100a2} of the
repository, every module compiles without \texttt{sorry}, and \texttt{lake build}
reports that \texttt{Erdos727.erdos727\_k2} depends only on the axioms
\texttt{propext}, \texttt{Classical.choice}, and \texttt{Quot.sound}; a
build-time assertion enforces the same for \texttt{erdos727\_k3}. This
was checked on the author's machine only; independent reproduction of the
build has not been recorded, and the argument presented here does not depend
on it. Independent expert review remains pending.

\section*{AI assistance and provenance}

This work was developed with assistance from \texttt{gpt-6-astra},
\texttt{fable 5.1}, and \texttt{gemini-3.8}. Responsibility for the
mathematical claims and the final manuscript rests with the author.

\begin{thebibliography}{9}
\bibitem{EGRS}
P.~Erd\H{o}s, R.~L.~Graham, I.~Z.~Ruzsa, and E.~G.~Straus,
\emph{On the prime factors of \(\binom{2n}{n}\)},
Mathematics of Computation \textbf{29} (1975), no.~129, 83--92.
The factorial-divisibility question appears on p.~90.
\href{https://www.renyi.hu/~p_erdos/1975-27.pdf}{Original article}.

\bibitem{DLMF}
NIST Digital Library of Mathematical Functions,
\emph{\S27.12: Asymptotic Formulas: Primes},
equations (27.12.5) and (27.12.8).
\url{https://dlmf.nist.gov/27.12}.
\end{thebibliography}
\end{document}
